\begin{longtblr}[
  caption = {UC 1.5: Manual Controlling},
]{
  colspec = {Q[l, 3.5cm] Q[l, 11cm]}, 
  hlines, vlines,
  row{1} = {bg=gray9, font=\bfseries}, 
  rowhead = 0, 
}
Field & Details \\
Use Case ID & 1.5 \\
Use Case Name & Manual Controlling \\
Actors & Primary: Parking Operators. \par Supporting: Gateway. \\
Description & Allows the Parking Operator to manually control the barrier gate via the software interface during emergencies, sensor failures, or non-standard entry/exit exceptions. \\
Pre-conditions & A hardware failure, unreadable plate, or emergency has occurred (Flow transferred from 1.1, 1.2, 1.3, or 1.4). The operator is logged into the system dashboard. \\
Post-conditions & The barrier state is manually changed (opened/closed), and the manual override event is securely logged for auditing. \\
Main Flow & 
1. The system alerts the Operator of a manual intervention request (e.g., license plate mismatch, hardware failure). \par
2. The Operator verifies the physical situation (e.g., visually checks the vehicle's actual license plate or identifies an emergency vehicle). \par
3. The Operator inputs the necessary missing data into the system (e.g., manual plate entry) or selects an override reason from a predefined list (e.g., ``Emergency'', ``System Error''). \par
4. The Operator clicks the manual override button on the system dashboard. \par
5. The system logs the manual override event, including the timestamp, Operator ID, and selected reason. \par
6. The system sends an execution command to the Gateway. \par
7. The Gateway physically triggers the barrier to open, allowing the vehicle to pass. \\
Exception Flows & 
\textbf{E1: Gateway Connection Lost:} The software cannot communicate with the physical Gateway due to a network or power failure. The system displays a connection error. The Operator must leave the booth and use a physical mechanical key to lift the barrier. \\
\end{longtblr}